\documentclass[10pt]{article}
\usepackage[margin=20mm]{geometry}
\usepackage{amsmath,amssymb,fontspec,unicode-math,xurl,hyperref}
\IfFontExistsTF{DejaVu Serif}{\setmainfont{DejaVu Serif}}{\setmainfont{Latin Modern Roman}}
\setmathfont{latinmodern-math.otf}
\setlength{\parindent}{0pt}\setlength{\parskip}{6pt}\sloppy
\begin{document}
Abstract: A smooth, pixel-valid radial decision boundary with a genuine bridge parameter disproves the identification of physical boundary curvature with the graph curvature of the auxiliary distance function. This invalidates a core geometric derivation. Separate findings concern a Taylor sign, a query-sum bound and three released-implementation discrepancies. The empirical two-thirds heuristic and published benchmark rankings are not refuted by these checks. Scientific adjudication remains pending.\par
Author: AI-assisted technical review\par
Date:\par
Title: Technical Verification of TtBA\par
\# Identity and inspected scope\par
Feiyang Wang, Xingquan Zuo, Hai Huang and Gang Chen, TtBA: Two-third Bridge Approach for Decision-Based Adversarial Attack, ICML 2025, PMLR 267, I25-43631 (local 0262), is reviewed in its 18-page proceedings version [Retained latex: \textbackslash{}cite\{paper\}]. Physical PDF and printed page numbers agree. The embedded appendices and relevant released code at commit `ea6d1649fccda0eb8392c9171e43cc9e0974acf9` were inspected [Retained latex: \textbackslash{}cite\{code\}]. This snapshot is not established as the publication-time executable.\par
The central geometric interpretation has a demonstrated invalid derivation. Separate local errors concern a Taylor coefficient, a query-count bound and the released evaluator. The empirical two-thirds rule remains an empirical choice; this review neither proves it optimal nor refutes its observed benchmark performance.\par
\# Definitions and the central geometric problem\par
For a correctly classified source image \(x\), let \(g(d)\) be the radial distance to the first adversarial boundary along \(\hat d=d/\|d\|_2\). Given current unit direction \(d\) and unit boundary normal \(N\), the paper defines \[u(k)=\frac{(1-k)d+kN}{\|(1-k)d+kN\|},\quad G(k)=g((1-k)d+kN).\] The actual boundary curve in image space is \(z(k)=x+G(k)u(k)\); it is not the Cartesian graph \((k,G(k))\).\par
Problem 1 (Curvature of the wrong geometric object: Major).  Appendix B.3, p. 13, Eq. (9), identifies the decision boundary’s curvature in the physical two-dimensional plane with \(b/(1+a^2)^{3/2}\), where \(a=G'(0)<0\), \(b=G''(0)>0\). This is the graph-curvature formula for \((k,G(k))\), not for \(z(k)\).\par
\#\# A counterexample with a valid bridge\par
Take a two-pixel image space \([0,1]^2\), source \(x=(1/2,1/2)\), and the smooth radial boundary about \(x\) \[\rho(\theta)=r[1-\sqrt3\sin\theta+v(1-\cos\theta)],
\qquad r=1/100,\quad v=22.\] Assign the original class to the interior and the adversarial class to the boundary and exterior. Its radial extrema are \(r(23\pm\sqrt{487})\), approximately \(0.00931924\) and \(0.450681\). Thus every radius is positive, the entire curve is inside the pixel square, and every ray from \(x\) has exactly one boundary intersection. The curve is smooth and regular because \(\rho^2+(\rho')^2>0\).\par
Equation reference: eq:G. At \(\theta=0\), the current boundary point is \(p=x+r(1,0)\), with \(d=(1,0)\). Since \(\rho'(0)=-\sqrt3 r\), the actual outward unit normal is \[N=\frac{(\rho(0),-\rho'(0))}{\sqrt{\rho(0)^2+\rho'(0)^2}}=(1/2,\sqrt3/2).\] Consequently the paper’s exact interpolation gives \[
G(k)=r\left[23-\frac{22-(19/2)k}{\sqrt{1-k+k^2}}\right].\] Direct differentiation yields \(a=-3r/2=-0.015\) and \(b=15r=0.15\), satisfying the stated derivative signs. Solving \(G(k)=G(0)=r\) gives the nonzero bridge \[k_{\rm bridge}=\frac{88}{525}=0.167619\ldots\in(0,1).\] Squaring is legitimate because \(22-(19/2)k\) is positive throughout \([0,1]\); the other root is \(k=0\). Moreover \(G(k)<r\) between these roots. This example therefore also has the bridge required by the method, rather than exploiting a missing-ray or missing-bridge pathology.\par
The actual plane-curve curvature magnitude at \(\theta=0\) is \[\kappa_{\rm physical}=
\frac{|\rho^2+2(\rho')^2-\rho\rho''|}
 {(\rho^2+(\rho')^2)^{3/2}}
=\frac{|7-v|}{8r}=187.5.\] Eq. (9) instead gives \(0.15/(1+0.015^2)^{3/2}=0.1499493892\). Both values concern the same boundary point and the same exact normal. The discrepancy is analytic; finite numerical differentiation is not needed.\par
\#\# Alternative reading and core dependence\par
If the intended object is the auxiliary graph \((k,G(k))\), Eq. (9) is valid for that graph. It then cannot be used as a derivation of the physical decision-boundary curvature. Page 3 identifies detecting varied boundary curvatures as the first contribution; pp. 5–6 use the bridge threshold to characterize narrow high-curvature regions and motivate the escape rule. The missing geometric correspondence therefore affects a central explanatory contribution and supports Major.\par
This witness tests the exact Eq. (9), not the separate higher-order-neglect step in Eq. (10). No assumption that its cubic remainder is negligible is used. The example does not show the attack fails or loses benchmark accuracy. The correct general plane-curve formula is \[\kappa=\frac{|\det(z',z'')|}{\|z'\|^3},\qquad
z'=G'u+Gu',\quad z''=G''u+2G'u'+Gu''.\] It restores the curvature calculation, but a full revised relation between this curvature and \(k_{\rm bridge}\) remains unproved. Even for graph curvature, an inverse proportionality inferred from \(-2a/k_{\rm bridge}\) requires controlling \(a\) across the comparison.\par
\# Independent local findings\par
Problem 2 (The cubic correction has the opposite sign: Minor).  Appendix B.4, p. 14, Eq. (19), expands the stationary point of \(a+bk+(c/2)k^2=0\) with a positive correction \(a^2c/(2b^3)\).\par
For \(b>0\) the branch continuous at \(c=0\) is \[k(c)=\frac{-b+\sqrt{b^2-2ac}}{c}
=-\frac ab-\frac{a^2c}{2b^3}+O(c^2).\] Substitution \(k=-a/b+tc\) gives \(bt+a^2/(2b^2)=0\), independently confirming the negative sign. With \(a=-0.1\), \(b=1\), \(c=0.02\), the exact root is \(0.0999001995\), the printed first-order expression is \(0.1001\), and the corrected expression is \(0.0999\). The root is inside \((0,1)\), the discriminant is positive, and \(G''=b+ck>0\). The quadratic relation \(k_{\min}=k_{\rm bridge}/2\) survives. The paper uses two-thirds empirically (p. 5 and Appendix G), so this local sign error is not inflated into a disproof of that heuristic.\par
Problem 3 (The claimed query upper bound is below the sum: Minor).  Appendix C, p. 15, Eqs. (21)–(22), sets \(q_i=24+\lceil30\sqrt i\rceil\) and claims \[Q_n=10+\sum_{i=1}^{n}q_i\leq24n+20n^{3/2}-34.\]\par
At \(n=1\), the exact left side is \(64\) and the claimed upper bound is \(10\). At \(n=57\) they are \(10115\) and \(9940.811256\), respectively. An increasing-function sum was replaced by an integral over \([1,n]\) in the wrong direction, and the ceiling contribution was dropped. For every integer \(n\geq1\), a valid conservative replacement is \[Q_n\leq10+25n+20[(n+1)^{3/2}-1],\] using \(\lceil t\rceil\leq t+1\) and \(\sum_{i=1}^n\sqrt i\leq\int_1^{n+1}\sqrt x\,dx\). In the paper’s modeled count, \(Q_{56}=9864\) and \(Q_{57}=10115\). The \(O(n^{3/2})\) order and qualitative small overhead survive. Actual boundary-search costs were stated approximately, so the empirical mean of 57 iterations is not refuted by this model-level correction.\par
Problem 4 (Released distortion is capped by the source-image norm: Minor).  Section 3 Eq. (1) and Section 5.2 use \(\|\tilde x-x\|_2\). The release, `TtBA.py`, lines 232–241, instead stores \(\min(\|x\|_2,\|\tilde x-x\|_2)\) in the distortion curve and uses it for the success threshold.\par
Let \(x=1/5\), classify \(z<4/5\) as the original class and \(z\geq4/5\) as adversarial, and set \(\epsilon=3/10\). The source is correctly classified, both source and boundary are legal pixels, and the true boundary distance is \(3/5>\epsilon\). The recorded value is \(1/5<\epsilon\), falsely satisfying the threshold. Clipping and the stated \(10^{-4}\) search tolerance do not erase this margin. Removing the cap and measuring the actual clamped adversarial perturbation corrects this metric definition, but no full release patch was executed. The source norm can exceed all encountered perturbations in other runs, so the cap may be inactive there. Publication-time logs are needed to determine historical impact; no published table is declared wrong from this snapshot alone.\par
Problem 5 (Normal-estimation settings differ from the stated recipe: Minor).  Algorithm 3, p. 16, uses normalized Gaussian directions with radius \(r_N=0.0003\). In `TtBA.py`, lines 117–122, the actual nonzero noise norm is \(\sigma\sqrt D/s\), and the full-dimensional branch samples uniform-cube rather than Gaussian noise.\par
Thus a 784-pixel input with \(s=1\) uses radius \(0.0084\), while a \(3\times224\times224\) input with \(s=4\) uses \(0.02909845\). These are direct algebraic consequences of the normalization, not estimated attack effects. Additionally, `dataset.py`, lines 255–262, resizes CIFAR100-ViT to width 224; the constructor’s width rule leaves its default \(s=4\), whereas p. 6 states \(s=1\) for datasets other than ImageNet. These changes may be intentional in the release, but correspondence with the published recipe is unresolved. No conclusion about which settings improve performance follows.\par
Problem 6 (Initialization fallback lacks a budget exit: Minor).  In `TtBA.py`, lines 271–287, once the initialization counter reaches \(Q/10\), the fallback tries all-white and all-black inputs without incrementing that counter or checking the total budget inside the loop.\par
If both fallback images retain the source class, the loop repeats them forever. These conditions are compatible with other adversarial images existing: in one dimension, take source \(x=0.2\) and an adversarial interval \((0.4,0.6)\). Both 0 and 1 are non-adversarial. A finite sequence of sufficiently negative Gaussian draws in the initial random phase misses the interval and has positive probability; after the fallback is reached, repeated 0/1 queries cannot escape it. This is a conditional release control-flow failure, not evidence that published runs followed that path. A bounded failure return or budget-aware restart is necessary; neither was implemented or benchmarked here.\par
\# Scope, counterchecks and outcome\par
The parent checked the original method, main experimental claims, Appendix B–G and the relevant released attack code; auxiliary static checks covered the data/model wrappers. Normal queries were found to be counted, with no established double counting. A completed iteration can cross a budget, but the summary code uses the preceding completed checkpoint for fixed-budget curves; this was not labeled a demonstrated budget-leakage error. The release also supplies radius scaling omitted from the literal Algorithm 2 call convention, so a pseudocode bracketing concern is not asserted as that particular implementation failure.\par
Original pages 13–15 were visually inspected to confirm the curvature expression, correction sign and inequality direction. Benchmark images, trained weights and datasets were not executed or regenerated; many figure panels were read only through captions and reported values. Four boundary-visualization PDFs in the source archive were preserved but not independently analyzed. Later revisions/errata were not surveyed; separate attachments remain unenumerated after acquisition-time HTTP 403.\par
Dimension | State\par
Assessment | Major geometric derivation; five Minor local/release findings\par
Adjudication | Pending; no human reviews\par
Repairs | Curvature identity, Taylor sign and sum bound analytically checked; no empirical/code repair\par
Author adoption | Not requested; unknown\par
Remaining work | Correct bridge–physical-curvature theory; historical source/run correspondence and empirical effects\par
9 F. Wang, X. Zuo, H. Huang and G. Chen. TtBA: Two-third Bridge Approach for Decision-Based Adversarial Attack. PMLR 267 (ICML 2025), 18-page proceedings PDF. https://proceedings.mlr.press/v267/wang25fg.html. OpenReview: `wEtNrpA8QH`. BUPTAIOC/TtBA, commit `ea6d1649fccda0eb8392c9171e43cc9e0974acf9`. https://github.com/BUPTAIOC/TtBA. Retrieval-time snapshot; static inspection only.\par
\# Reproducibility and roles\par
Original PDF SHA-256: 826799647b7967c40c919ece2665b366c8b01891d03eba77ccc1ff2d669095e3. Under `verification/0262/`, `check.py` and `check.json` contain the parent’s independent finite calculations; `theory\_notes.json` and `code\_static.json` retain native auxiliary cross-checks and alternatives. The radial bridge counterexample originated with the parent and was independently checked. Original-page images are `s13.png`–`s16.png`. All source/archive/member identities remain in the unchanged source manifest and the actual byte-check records.\par
The parent owns scientific integration, independent calculations, exact-TeX compilation and final rendered-PDF inspection. Source-integrity, theory and static-source auxiliary roles are recorded separately. The exposed system label is GPT-6 Astra Pro; requested model is 6 Pro; unobserved UI label and conversation URL remain null. There was no author-code execution, software installation, empirical reproduction, author contact or human adjudication. Actual build commands, repaired compile failures if any, artifact hashes and visual records are preserved with this initial-review report.\par
\end{document}
